\section{Related Work}

Our work synthesizes three research threads --- sub-quadratic architectures, efficient adaptation, and loss landscape analysis --- to address the unexplored question of task-dependent adaptation under architecture conversion.

\subsection{Sub-Quadratic Architectures}

The Transformer's quadratic attention complexity has motivated extensive work on sub-quadratic alternatives. Linear attention methods (Performer, cosFormer) replace $\text{softmax}(QK^T)V$ with kernel-based approximations, achieving $O(n)$ complexity but often degrading on tasks requiring precise attention patterns. State space models, particularly S4 and its successor Mamba, offer a fundamentally different approach: selective scan operations that evolve hidden state sequentially, achieving transformer-quality performance with linear complexity.

Mamba demonstrates SSM-attention duality, showing that attention and state space operations are structurally connected. Mamba-2 formalizes this duality, enabling principled conversion between architectures. However, existing evaluations focus on pretraining quality rather than post-conversion adaptation behavior. Hybrid architectures like Jamba interleave Mamba layers with sparse attention, suggesting that different architectural components serve different functions --- a perspective our work makes precise through task-dependent analysis.

\subsection{Efficient Adaptation Methods}

LoRA introduced low-rank weight decomposition for efficient fine-tuning, enabling adaptation of large models by training only $O(\text{rank} \times d)$ parameters. The method's success on Transformers led to variants (QLoRA, DoRA) and adoption across architectures. Yet all LoRA studies assume fixed architecture, missing how architecture change affects adaptation efficiency. Our work fills this gap by characterizing LoRA behavior under Transformer-to-Mamba conversion.

\subsection{Loss Landscape Analysis}

Sharpness-aware minimization (SAM) established the connection between loss landscape geometry and generalization: flatter minima generalize better. We extend landscape analysis to architecture conversion, discovering that conversion itself transforms landscape geometry (219\% sharpness change). More importantly, we connect landscape sharpness to LoRA effective rank ($\rho=1.0$), providing mechanistic explanation for task-dependent adaptation efficiency.
